\textbf{Phần 1. Trắc nghiệm lựa chọn}
\setcounterpageref{ex}{0}

\begin{ex} %[4D8N41]
Khi nhân một số tự nhiên với 10, ta chỉ việc
\choice
{Thêm hai chữ số 0 vào bên phải số đó}
{\True Thêm một chữ số 0 vào bên phải số đó}
{Thêm ba chữ số 0 vào bên phải số đó}
{Nhân số đó với 1}
\loigiai{
Nhân với 10 viết thêm một chữ số 0 vào bên phải.\\
\textbf{Đáp án của câu hỏi là B}
}
\end{ex}

\begin{ex} %[4D8N41]
Khi chia một số tròn trăm cho 100, ta chỉ việc
\choice
{Bớt một chữ số 0 ở bên phải số đó}
{Bớt ba chữ số 0 ở bên phải số đó}
{\True Bớt hai chữ số 0 ở bên phải số đó}
{Không thay đổi}
\loigiai{
Chia cho 100 bớt đi hai chữ số 0 ở bên phải.\\
\textbf{Đáp án của câu hỏi là C}
}
\end{ex}

\begin{ex} %[4D8H41]
Kết quả của phép nhân $345 \times 100$ là
\choice
{$3450$}
{$345\,000$}
{$3450$}
{\True $34\,500$}
\loigiai{
$345 \times 100 = 34\,500$.\\
\textbf{Đáp án của câu hỏi là D}
}
\end{ex}

\begin{ex} %[4D8H41]
Kết quả của phép chia $280\,000 : 1000$ là
\choice
{\True $280$}
{$28$}
{$2800$}
{$28\,000$}
\loigiai{
$280\,000 : 1000 = 280$.\\
\textbf{Đáp án của câu hỏi là A}
}
\end{ex}

\begin{ex} %[4D8H41]
Điền số thích hợp vào chỗ chấm: $450 \times \dots = 45\,000$.
\choice
{$10$}
{$1000$}
{$10\,000$}
{\True $100$}
\loigiai{
$45\,000 : 450 = 100$.\\
\textbf{Đáp án của câu hỏi là D}
}
\end{ex}

\begin{ex} %[4D8H41]
Điền số thích hợp vào chỗ chấm: $670\,000 : \dots = 670$.
\choice
{\True $1000$}
{$100$}
{$10$}
{$10\,000$}
\loigiai{
$670\,000 : 670 = 1000$.\\
\textbf{Đáp án của câu hỏi là A}
}
\end{ex}

\begin{ex} %[4D8V41]
Một bao gạo nặng $50\text{ kg}$. Hỏi $100$ bao gạo như thế nặng bao nhiêu tấn?
\choice
{$0{,}5\text{ tấn}$}
{$50\text{ tấn}$}
{\True $5\text{ tấn}$}
{$500\text{ tấn}$}
\loigiai{
$50 \times 100 = 5000\text{ kg} = 5\text{ tấn}$.\\
\textbf{Đáp án của câu hỏi là C}
}
\end{ex}

\begin{ex} %[4D8V41]
Một xí nghiệp đóng $150\,000$ gói kẹo vào các thùng, mỗi thùng $100$ gói. Số thùng kẹo đóng được là
\choice
{$150$ thùng}
{\True $1500$ thùng}
{$15\,000$ thùng}
{$15$ thùng}
\loigiai{
$150\,000 : 100 = 1500$ thùng.\\
\textbf{Đáp án của câu hỏi là B}
}
\end{ex}

\begin{ex} %[4D8V41]
Tính giá trị biểu thức: $125 \times 100 + 450 \times 10$.
\choice
{$16\,000$}
{$18\,000$}
{$17\,500$}
{\True $17\,000$}
\loigiai{
$12\,500 + 4500 = 17\,000$.\\
\textbf{Đáp án của câu hỏi là D}
}
\end{ex}

\begin{ex} %[4D8V41]
Tìm $x$ biết: $x \times 100 = 24\,000 : 10$.
\choice
{$240$}
{\True $24$}
{$2400$}
{$2{,}4$}
\loigiai{
$x \times 100 = 2400 \Rightarrow x = 24$.\\
\textbf{Đáp án của câu hỏi là B}
}
\end{ex}

\begin{ex} %[4D8V41]
Một mảnh đất hình chữ nhật có chiều dài $40\text{ m}$, chiều rộng $25\text{ m}$. Diện tích mảnh đất theo mét vuông là
\choice
{\True $1000\text{ m}^2$}
{$100\text{ m}^2$}
{$10\,000\text{ m}^2$}
{$65\text{ m}^2$}
\loigiai{
$40 \times 25 = 1000\text{ m}^2$.\\
\textbf{Đáp án của câu hỏi là A}
}
\end{ex}

\begin{ex} %[4D8C41]
Nếu gấp một số lên 100 lần rồi giảm đi 10 lần thì được số $450$. Số ban đầu là
\choice
{$4500$}
{$450$}
{\True $45$}
{$4{,}5$}
\loigiai{
Gấp 100 lần rồi giảm 10 lần tức là gấp 10 lần. Số ban đầu là $450 : 10 = 45$.\\
\textbf{Đáp án của câu hỏi là C}
}
\end{ex}

\textbf{Phần 2. Trắc nghiệm đúng sai}
\setcounterpageref{ex}{0}

\begin{ex} %[4D8H41]
Xét tính đúng sai của các phép nhân, chia:
\choiceTF
{\True $78 \times 1000 = 78\,000$.}
{$450\,000 : 100 = 450$.}
{\True $450\,000 : 100 = 4500$.}
{\True $120 \times 10 = 1200$.}
\loigiai{
\begin{itemize}
    \item $78 \times 1000 = 78\,000$. Đúng.
    \item $450\,000 : 100 = 4500 \neq 450$. Ý 2 Sai, Ý 3 Đúng.
    \item $120 \times 10 = 1200$. Đúng.
\end{itemize}
\begin{itemize}
    \item a) \textbf{Đúng}.
    \item b) \textbf{Sai}.
    \item c) \textbf{Đúng}.
    \item d) \textbf{Đúng}.
\end{itemize}
}
\end{ex}

\begin{ex} %[4D8V41]
Một trang trại nuôi bò sữa mỗi ngày vắt được $1200\text{ l}$ sữa. Sữa được đóng vào các chai $1\text{ l}$ rồi đóng vào các thùng, mỗi thùng $10$ chai. Xét tính đúng sai:
\choiceTF
{\True Mỗi ngày trang trại đóng được $1200$ chai sữa.}
{\True Mỗi ngày trang trại đóng được $120$ thùng sữa.}
{Trong 10 ngày trang trại đóng được $1000$ thùng sữa.}
{\True Trong 10 ngày trang trại đóng được $1200$ thùng sữa.}
\loigiai{
\begin{itemize}
    \item 1200 chai: Đúng.
    \item $1200 : 10 = 120$ thùng: Đúng.
    \item 10 ngày: $120 \times 10 = 1200$ thùng $\neq 1000$ thùng. Ý 3 Sai, Ý 4 Đúng.
\end{itemize}
\begin{itemize}
    \item a) \textbf{Đúng}.
    \item b) \textbf{Đúng}.
    \item c) \textbf{Sai}.
    \item d) \textbf{Đúng}.
\end{itemize}
}
\end{ex}

\textbf{Phần 3. Trả lời ngắn}
\setcounterpageref{ex}{0}

\begin{ex} %[4D8H41]
Tính tích: $560 \times 100$.
\shortans{56000}
\loigiai{
\textbf{Đáp số:} $56000$\\
$56\,000$.
}
\end{ex}

\begin{ex} %[4D8H41]
Tính thương: $890\,000 : 1000$.
\shortans{890}
\loigiai{
\textbf{Đáp số:} $890$\\
$890$.
}
\end{ex}

\begin{ex} %[4D8V41]
Tìm $x$ biết: $x \times 100 = 75\,000$.
\shortans{750}
\loigiai{
\textbf{Đáp số:} $750$\\
$x = 750$.
}
\end{ex}

\begin{ex} %[4D8V41]
Một trường tiểu học mua 100 quyển sách tham khảo giá $35\,000$ đồng một quyển. Tổng số tiền mua sách là bao nhiêu nghìn đồng (chỉ ghi số nghìn đồng)?
\shortans{3500}
\loigiai{
\textbf{Đáp số:} $3500$\\
$35\,000 \times 100 = 3\,500\,000$ đồng $= 3500$ nghìn đồng.
}
\end{ex}

\begin{ex} %[4D8V41]
Tính giá trị biểu thức: $450 \times 100 - 35\,000$.
\shortans{10000}
\loigiai{
\textbf{Đáp số:} $10000$\\
$45\,000 - 35\,000 = 10\,000$.
}
\end{ex}

\begin{ex} %[4D8C41]
Một số tròn nghìn chia cho 1000 được số bé nhất có 3 chữ số khác nhau. Tìm số tròn nghìn đó.
\shortans{102000}
\loigiai{
\textbf{Đáp số:} $102000$\\
Số bé nhất có 3 chữ số khác nhau là 102. Số cần tìm: $102 \times 1000 = 102\,000$.
}
\end{ex}
